\chapter{Background}\label{ch:background}

This chapter introduces the concepts the rest of the document depends on. It
is deliberately neutral: it presents established ideas --- separation of
concerns in parallel programming, the dimensions of IO semantics, static versus
dynamic scheduling, Petri nets, and differential testing --- without yet
assuming the contribution. A reader already fluent in any of these sections may
skip it.

\section{Parallel decomposition and separation of concerns}\label{sec:bg-decomp}

A sequential program states \emph{what} to compute. A parallel program must also
state \emph{how} the work is divided across processing elements and \emph{how}
those elements coordinate. These are different concerns, and conflating them is
the source of much of the difficulty in parallel programming: a change to the
data layout for one machine forces edits all through code that expresses the
mathematics, and the mathematics becomes hard to read through the layout.

The decomposition itself has two axes. \emph{Spatial} decomposition divides data
or iterations across workers that run concurrently --- domain decomposition of a
grid, blocking of a matrix, partitioning of an array. \emph{Temporal}
decomposition orders work within a worker and overlaps it across workers ---
loop tiling, software pipelining (overlapping successive loop iterations so
workers stay busy), and double buffering (alternating between two buffers so a
producer can fill one while the consumer drains the other). Both axes are
scheduling decisions: they change performance and resource use, but, applied
correctly, they do not change the computed result.

The idea that the result-defining part of a program should be written once and
the performance-defining part chosen separately has a clear lineage in
high-performance compute. Halide~\cite{ragan-kelley2013halide} crystallised it
as the explicit split between an \emph{algorithm} (pure functions over an index
domain) and a \emph{schedule} (the tiling, fusion, vectorisation, and
parallelism applied to it), and showed that a single algorithm could be
retargeted across schedules by orders of magnitude in performance without
touching its definition. The polyhedral
compilers~\cite{baghdadi2019tiramisu,bondhugula2008pluto} pursue the same
separation through a more powerful mathematical model. What unites them is the
discipline that the algorithm never mentions the schedule. This document
adopts that discipline and extends it from scheduling computation \emph{within a
single machine} (intra-node) to decomposition and IO \emph{across many machines
or workers} (inter-node). A middle ground between fully manual and fully
automatic decomposition is to have the programmer annotate the decomposition
explicitly --- in a separate schedule sublanguage --- while leaving the compiler
to infer the resulting transfers from those annotations; that is the approach
taken here and developed in \cref{ch:language}.

\section{IO semantics as a programming concern}\label{sec:bg-io}

When two workers exchange data, the exchange is governed by choices that are
usually made implicitly by whatever runtime is in use, but which materially
affect correctness and performance. These choices form a small set of orthogonal
dimensions:

\begin{description}[leftmargin=*]
  \item[Synchronous vs.\ asynchronous.] Does the producer block until the
        transfer completes, or does it initiate the transfer and continue,
        completing it later? Asynchrony enables latency hiding but requires
        somewhere to track the in-flight operation.
  \item[Buffered vs.\ unbuffered.] Is there storage that decouples producer and
        consumer rates, allowing the producer to run ahead by a bounded number
        of items, or must each transfer rendezvous directly?
  \item[Event-driven vs.\ polled.] Is the consumer woken when data arrives (by
        an interrupt, a condition variable, or an event-loop readiness
        notification), or does it repeatedly test for arrival?
  \item[Transport.] Do the bytes move through shared memory, a loopback or
        network socket, a Unix domain socket, a cluster interconnect via a
        message-passing library, or a DMA engine on an embedded device?
\end{description}

These dimensions are largely independent of one another and of the algorithm. A
shared-memory threads runtime such as OpenMP~\cite{dagum1998openmp} fixes one
combination; a message-passing library such as MPI~\cite{mpiforum2021} offers
several (blocking and non-blocking sends, buffered and synchronous modes) but
within its own model; an embedded async framework such as
Embassy~\cite{embassy} fixes yet another around interrupts and DMA. Crucially,
none of them lets these choices be \emph{stated} as part of a schedule and
\emph{swapped} without rewriting the communication code. That gap --- IO
semantics as an unspoken property of the chosen runtime rather than a
first-class, selectable directive --- is the gap this work targets.

\section{Static versus dynamic scheduling}\label{sec:bg-scheduling}

A \emph{schedule} assigns operations to workers and orders them in time.
Scheduling may happen dynamically, at run time --- as in work-stealing thread
pools or dataflow runtimes that fire operations when their inputs become
available --- or statically, at compile time, fixing the order before the
program runs.

Dynamic scheduling adapts to data-dependent and machine-dependent variation, at
the cost of runtime overhead, nondeterminism, and analyses that can only be
approximate. Static scheduling gives up that adaptivity but gains three
properties this work depends on. First, \emph{determinism}: a fixed firing order
means a fixed result and a fixed communication pattern, which is what makes
byte-identical cross-backend comparison meaningful. Second,
\emph{analysability}: with the order known at compile time, properties such as
buffer sufficiency and freedom from deadlock become decidable by replaying that
one order, rather than searching a reachable state space. Third, \emph{suitability
for embedded real-time targets}, where dynamic scheduling is often unacceptable
regardless. Static scheduling of dataflow has a long pedigree --- synchronous
dataflow~\cite{lee1987sdf} schedules a restricted graph entirely at compile time
precisely to remove runtime overhead --- and this work sits firmly in that
static tradition. The scope restrictions stated in \cref{sec:intro-scope} ---
affine bounds, single-assignment arrays, and a statically fixed decomposition ---
are exactly what place this work in the static-scheduling camp and make the
analyses of the following sections tractable.

\section{Petri nets}\label{sec:bg-petri}

A Petri net~\cite{petri1962kommunikation,murata1989petri} is a bipartite
directed graph with two kinds of node: \emph{places} and \emph{transitions}.
Places hold \emph{tokens}; a marking is an assignment of a token count to every
place. A transition is \emph{enabled} when each of its input places holds enough
tokens, and \emph{firing} it removes tokens from its input places and deposits
tokens on its output places, producing a new marking. The net's behaviour is the
set of markings reachable from an initial marking by firing enabled transitions.

A handful of standard properties characterise such a net. A net is
\emph{bounded} if no place can ever exceed a stated token capacity, and
\emph{$k$-bounded} if that capacity is $k$. It is \emph{live} if no transition is
permanently disabled. A \emph{deadlock} is a reachable marking in which no
transition is enabled --- the system cannot make progress. For general Petri
nets these properties are decidable but expensive; for restricted subclasses
they become cheap.

Petri nets are a natural model for the problem at hand because their vocabulary
maps directly onto statically scheduled communication. A transition models a
kernel firing, a data transfer, or a synchronisation; a place models a data
slot, a communication channel, or a barrier; a token models the presence of a
value or a unit of buffer credit; a place's capacity models a declared buffer
size; and tokens placed on a channel in the initial marking model the
head-start that double-buffering or software-pipelining provides. Under this
mapping, ``does this schedule's communication stay within its declared
buffers?'' is exactly boundedness, and ``can this schedule's communication
hang?'' is exactly the reachability of a deadlocked marking.
\Cref{ch:architecture} shows how restricting the net to a statically determined
firing order collapses these checks to a single deterministic replay.

\section{Determinism and differential testing}\label{sec:bg-determinism}

\emph{Differential testing}~\cite{mckeeman1998differential} establishes
confidence in a system by comparing the outputs of multiple implementations that
are supposed to agree: where they differ, at least one is wrong, and the
disagreement localises the fault. It is especially powerful for compilers, where
an oracle for ``the correct output'' is otherwise hard to obtain; the
Csmith~\cite{yang2011csmith} work famously used differential testing across C
compilers to find hundreds of bugs by generating programs whose output every
correct compiler must agree on.

Differential testing requires a deterministic notion of agreement. Two
implementations can only be compared byte-for-byte if both are deterministic for
the same input. This is why the algorithm class in this work is restricted to
deterministic arithmetic (\cref{sec:intro-scope}): integer computation, or
floating-point computation whose reductions do not reorder, guarantees that a
single reference output exists for every example. Given that determinism,
differential testing becomes a means not only of finding bugs but of
\emph{falsifying a design claim}. If one algorithm is compiled under many
schedules and many backends and all outputs are byte-identical to a single
hand-written reference, then the claim that the algorithm is independent of
schedule and backend has survived a concrete attempt to break it; a single
differing byte refutes it. \Cref{ch:validation} develops this rig in full.
